% Options for packages loaded elsewhere
% Options for packages loaded elsewhere
\PassOptionsToPackage{unicode}{hyperref}
\PassOptionsToPackage{hyphens}{url}
\PassOptionsToPackage{dvipsnames,svgnames,x11names}{xcolor}
%
\documentclass[
  letterpaper,
  DIV=11,
  numbers=noendperiod]{scrartcl}
\usepackage{xcolor}
\usepackage{amsmath,amssymb}
\setcounter{secnumdepth}{5}
\usepackage{iftex}
\ifPDFTeX
  \usepackage[T1]{fontenc}
  \usepackage[utf8]{inputenc}
  \usepackage{textcomp} % provide euro and other symbols
\else % if luatex or xetex
  \usepackage{unicode-math} % this also loads fontspec
  \defaultfontfeatures{Scale=MatchLowercase}
  \defaultfontfeatures[\rmfamily]{Ligatures=TeX,Scale=1}
\fi
\usepackage{lmodern}
\ifPDFTeX\else
  % xetex/luatex font selection
\fi
% Use upquote if available, for straight quotes in verbatim environments
\IfFileExists{upquote.sty}{\usepackage{upquote}}{}
\IfFileExists{microtype.sty}{% use microtype if available
  \usepackage[]{microtype}
  \UseMicrotypeSet[protrusion]{basicmath} % disable protrusion for tt fonts
}{}
\makeatletter
\@ifundefined{KOMAClassName}{% if non-KOMA class
  \IfFileExists{parskip.sty}{%
    \usepackage{parskip}
  }{% else
    \setlength{\parindent}{0pt}
    \setlength{\parskip}{6pt plus 2pt minus 1pt}}
}{% if KOMA class
  \KOMAoptions{parskip=half}}
\makeatother
% Make \paragraph and \subparagraph free-standing
\makeatletter
\ifx\paragraph\undefined\else
  \let\oldparagraph\paragraph
  \renewcommand{\paragraph}{
    \@ifstar
      \xxxParagraphStar
      \xxxParagraphNoStar
  }
  \newcommand{\xxxParagraphStar}[1]{\oldparagraph*{#1}\mbox{}}
  \newcommand{\xxxParagraphNoStar}[1]{\oldparagraph{#1}\mbox{}}
\fi
\ifx\subparagraph\undefined\else
  \let\oldsubparagraph\subparagraph
  \renewcommand{\subparagraph}{
    \@ifstar
      \xxxSubParagraphStar
      \xxxSubParagraphNoStar
  }
  \newcommand{\xxxSubParagraphStar}[1]{\oldsubparagraph*{#1}\mbox{}}
  \newcommand{\xxxSubParagraphNoStar}[1]{\oldsubparagraph{#1}\mbox{}}
\fi
\makeatother


\usepackage{longtable,booktabs,array}
\usepackage{calc} % for calculating minipage widths
% Correct order of tables after \paragraph or \subparagraph
\usepackage{etoolbox}
\makeatletter
\patchcmd\longtable{\par}{\if@noskipsec\mbox{}\fi\par}{}{}
\makeatother
% Allow footnotes in longtable head/foot
\IfFileExists{footnotehyper.sty}{\usepackage{footnotehyper}}{\usepackage{footnote}}
\makesavenoteenv{longtable}
\usepackage{graphicx}
\makeatletter
\newsavebox\pandoc@box
\newcommand*\pandocbounded[1]{% scales image to fit in text height/width
  \sbox\pandoc@box{#1}%
  \Gscale@div\@tempa{\textheight}{\dimexpr\ht\pandoc@box+\dp\pandoc@box\relax}%
  \Gscale@div\@tempb{\linewidth}{\wd\pandoc@box}%
  \ifdim\@tempb\p@<\@tempa\p@\let\@tempa\@tempb\fi% select the smaller of both
  \ifdim\@tempa\p@<\p@\scalebox{\@tempa}{\usebox\pandoc@box}%
  \else\usebox{\pandoc@box}%
  \fi%
}
% Set default figure placement to htbp
\def\fps@figure{htbp}
\makeatother


% definitions for citeproc citations
\NewDocumentCommand\citeproctext{}{}
\NewDocumentCommand\citeproc{mm}{%
  \begingroup\def\citeproctext{#2}\cite{#1}\endgroup}
\makeatletter
 % allow citations to break across lines
 \let\@cite@ofmt\@firstofone
 % avoid brackets around text for \cite:
 \def\@biblabel#1{}
 \def\@cite#1#2{{#1\if@tempswa , #2\fi}}
\makeatother
\newlength{\cslhangindent}
\setlength{\cslhangindent}{1.5em}
\newlength{\csllabelwidth}
\setlength{\csllabelwidth}{3em}
\newenvironment{CSLReferences}[2] % #1 hanging-indent, #2 entry-spacing
 {\begin{list}{}{%
  \setlength{\itemindent}{0pt}
  \setlength{\leftmargin}{0pt}
  \setlength{\parsep}{0pt}
  % turn on hanging indent if param 1 is 1
  \ifodd #1
   \setlength{\leftmargin}{\cslhangindent}
   \setlength{\itemindent}{-1\cslhangindent}
  \fi
  % set entry spacing
  \setlength{\itemsep}{#2\baselineskip}}}
 {\end{list}}
\usepackage{calc}
\newcommand{\CSLBlock}[1]{\hfill\break\parbox[t]{\linewidth}{\strut\ignorespaces#1\strut}}
\newcommand{\CSLLeftMargin}[1]{\parbox[t]{\csllabelwidth}{\strut#1\strut}}
\newcommand{\CSLRightInline}[1]{\parbox[t]{\linewidth - \csllabelwidth}{\strut#1\strut}}
\newcommand{\CSLIndent}[1]{\hspace{\cslhangindent}#1}



\setlength{\emergencystretch}{3em} % prevent overfull lines

\providecommand{\tightlist}{%
  \setlength{\itemsep}{0pt}\setlength{\parskip}{0pt}}



 


\KOMAoption{captions}{tableheading}
\makeatletter
\@ifpackageloaded{caption}{}{\usepackage{caption}}
\AtBeginDocument{%
\ifdefined\contentsname
  \renewcommand*\contentsname{Table of contents}
\else
  \newcommand\contentsname{Table of contents}
\fi
\ifdefined\listfigurename
  \renewcommand*\listfigurename{List of Figures}
\else
  \newcommand\listfigurename{List of Figures}
\fi
\ifdefined\listtablename
  \renewcommand*\listtablename{List of Tables}
\else
  \newcommand\listtablename{List of Tables}
\fi
\ifdefined\figurename
  \renewcommand*\figurename{Figure}
\else
  \newcommand\figurename{Figure}
\fi
\ifdefined\tablename
  \renewcommand*\tablename{Table}
\else
  \newcommand\tablename{Table}
\fi
}
\@ifpackageloaded{float}{}{\usepackage{float}}
\floatstyle{ruled}
\@ifundefined{c@chapter}{\newfloat{codelisting}{h}{lop}}{\newfloat{codelisting}{h}{lop}[chapter]}
\floatname{codelisting}{Listing}
\newcommand*\listoflistings{\listof{codelisting}{List of Listings}}
\makeatother
\makeatletter
\usepackage{pdflscape}
\makeatother
\makeatletter
\makeatother
\makeatletter
\@ifpackageloaded{caption}{}{\usepackage{caption}}
\@ifpackageloaded{subcaption}{}{\usepackage{subcaption}}
\makeatother
\usepackage{bookmark}
\IfFileExists{xurl.sty}{\usepackage{xurl}}{} % add URL line breaks if available
\urlstyle{same}
\hypersetup{
  pdftitle={Comparing Habitat Suitability Modeling Approaches for the Eastern Oyster in the Pensacola Bay System},
  pdfauthor={M. Fazio},
  colorlinks=true,
  linkcolor={blue},
  filecolor={Maroon},
  citecolor={Blue},
  urlcolor={Blue},
  pdfcreator={LaTeX via pandoc}}


\title{Comparing Habitat Suitability Modeling Approaches for the Eastern
Oyster in the Pensacola Bay System}
\author{M. Fazio}
\date{10/20/2026}
\begin{document}
\maketitle

\renewcommand*\contentsname{Table of contents}
{
\hypersetup{linkcolor=}
\setcounter{tocdepth}{3}
\tableofcontents
}

\section{Introduction}\label{introduction}

The eastern oyster (\emph{Crassostrea virginica}) are viewed as a
foundational species as they are directly responsible for creating and
providing numerous ecosystem goods and services such as food and habitat
production, shoreline stabilization and protection, water quality
improvement, nutrient cycling, and carbon sequestration (Geselbracht et
al. 2024; Grabowski et al. 2012; Barnes et al. 2007). Additionally, they
provide numerous economic benefits such as supporting oyster fisheries,
recreational harvesting, tourism, along with the associated employment
that goes along with those industries (Geselbracht et al. 2024; Peyre,
Marshall, and Sable 2021). In one study, Grabowski et al. (2012)
estimated the average annual value to range from \$10,325 to \$99,421
per hectare (2011 Dollars).

In recent decades, oyster reefs have experienced substantial declines,
with only an estimated 15\% remaining worldwide, and 50 to 80\% being
lost within the Gulf of (America) alone (Pollack et al. 2012; Peyre,
Marshall, and Sable 2021). The decline is largely attributed to a
combination of anthropogenic activities such us over harvesting, habitat
alteration, sedimentation and water quality decline, disease and
changing hydrology (Peyre, Marshall, and Sable 2021; Geselbracht et al.
2024). In response to these sharp declines, many restoration efforts
have been undertaken with varying strategies and success rates;
recently, planning documents released by the Deepwater Horizon Natural
resource Damage Assessment Trustees called for large-scale development
of reef networks to increase resilience (Geselbracht et al. 2024).

\subsection{Oyster decline in the Pensacola Bay
System}\label{oyster-decline-in-the-pensacola-bay-system}

On such system to experience oyster decline is the Pensacola Bay System
(PBS). The PBS is a large estuarine system located in Escambia and Santa
Rosa counties in northwest Florida. The system is comprised of multiple
waterbody's including Pensacola Bay, Escambia Bay, Blackwater Bay, East
Bay, and Santa Rosa Sound; and is influenced heavily by freshwater input
from the Escambia, Blackwater and Yellow Rivers (Christensen et al.
1997). Historically, the PBS supported extensive oyster reefs that
provided critical habitat for a variety of species and substantial
economic value to the surrounding communities (Christensen et al. 1997).
However, over the past several decades, the extent of oyster habitat
within the PBS has declined significantly due to a combination of
factors including over harvesting, habitat degradation, and water
quality issues. With some estimates indicating a decline of 15,000ha to
just 91.5 ha of oyster habitat between the years of 1883 and 2021
(Geselbracht et al. 2024). This decline has prompted efforts to restore
oyster reefs within the PBS (Pensacola and Perdido Bays Estuary Program
n.d.; The Nature Conservancy 2021), with multiple studies aiming to
identify the most suitable areas within the estuary to focus restoration
efforts (Geselbracht et al. 2024; Christensen et al. 1997; Pensacola and
Perdido Bays Estuary Program n.d.).

With limited resources available for restoration, it is critical to
identify areas where chances of oyster survival and long-term reef
establishment are at their highest, and where poor site selection could
ultimately lead to unsatisfactory results. As with any ecosystem, the
environmental conditions within the PBS vary greatly in time and space,
and such variations can have great impacts on oyster survival and
establishment. Therefore, it is important to understand the
environmental factors that drive oyster occurrence so that scientist and
policy makers can make informed decisions regarding said restoration
site selection (Pollack et al. 2012; Barnes et al. 2007; Geselbracht et
al. 2024).

\subsection{Environmental controls on Eastern oyster
habitat}\label{environmental-controls-on-eastern-oyster-habitat}

Oyster recruitment, survival, and reef establishment are influenced by
complex interactions between numerous physical, chemical, and biological
environmental factors. As such, determining the most influential factors
is a multivariate problem that requires careful consideration of each
factor and its interactions with others (Peyre, Marshall, and Sable
2021; Pollack et al. 2012; Barnes et al. 2007; Christensen et al. 1997).
To help narrow the focus, the most commonly cited environmental factors
were examined in the literature and are summarized here.

Perhaps the most influential and widely reported environmental driver of
oyster suitability is salinity. Representing an important environmental
gradient within estuarine systems that influences physiology,
recruitment, disease, predation, and survival. Other important
characteristics often reported can be seen in
\textbf{?@tbl-environmental-drivers} and include: dissolved oxygen,
temperature, depth, turbidity, and bottom type (substrate) (Peyre,
Marshall, and Sable 2021; Pollack et al. 2012; Barnes et al. 2007;
Geselbracht et al. 2024; Christensen et al. 1997).

\begin{landscape}

\begin{longtable}[]{@{}
  >{\raggedright\arraybackslash}p{(\linewidth - 4\tabcolsep) * \real{0.1236}}
  >{\raggedright\arraybackslash}p{(\linewidth - 4\tabcolsep) * \real{0.5491}}
  >{\raggedright\arraybackslash}p{(\linewidth - 4\tabcolsep) * \real{0.3273}}@{}}
\caption{Common environmental drivers of Eastern Oyster habitat
suitability.}\tabularnewline
\toprule\noalign{}
\begin{minipage}[b]{\linewidth}\raggedright
Factor
\end{minipage} & \begin{minipage}[b]{\linewidth}\raggedright
Description
\end{minipage} & \begin{minipage}[b]{\linewidth}\raggedright
Sources
\end{minipage} \\
\midrule\noalign{}
\endfirsthead
\toprule\noalign{}
\begin{minipage}[b]{\linewidth}\raggedright
Factor
\end{minipage} & \begin{minipage}[b]{\linewidth}\raggedright
Description
\end{minipage} & \begin{minipage}[b]{\linewidth}\raggedright
Sources
\end{minipage} \\
\midrule\noalign{}
\endhead
\bottomrule\noalign{}
\endlastfoot
Salinity & Controls feeding, growth, reproduction, recruitment, disease,
predation, and survival. & (Barnes et al. 2007; Birch et al. 2021;
Christensen et al. 1997; Geselbracht et al. 2024; Peyre, Marshall, and
Sable 2021; Pollack et al. 2012) \\
Temperature & Regulates metabolism, filtration, growth, spawning, and
larval development. & (Barnes et al. 2007; Christensen et al. 1997;
Geselbracht et al. 2024; Peyre, Marshall, and Sable 2021; Pollack et al.
2012) \\
Substrate type & Stable hard substrate supports larval attachment and
reef persistence, whereas soft or rapidly accumulating sediments can
bury or destabilize oysters. & (Barnes et al. 2007; Birch et al. 2021;
Christensen et al. 1997; Geselbracht et al. 2024; Peyre, Marshall, and
Sable 2021) \\
Dissolved oxygen & Adequate oxygen supports respiration, feeding, and
growth. & (Christensen et al. 1997; Geselbracht et al. 2024; Peyre,
Marshall, and Sable 2021; Pollack et al. 2012) \\
Hydrology / water movement & Water movement regulates salinity and
transports larvae, oxygen, and food. & (Barnes et al. 2007; Geselbracht
et al. 2024; Peyre, Marshall, and Sable 2021) \\
Depth / bathymetry & Depth influences exposure, temperature, oxygen
conditions, and restoration feasibility. & (Christensen et al. 1997;
Geselbracht et al. 2024; Pollack et al. 2012) \\
Turbidity / sedimentation & High suspended-sediment loads can interfere
with filtration and settlement. Sediment deposition can smoother oyster
reefs. & (Peyre, Marshall, and Sable 2021; Pollack et al. 2012) \\
Food availability / chlorophyll a & Chlorophyll a is often used as a
proxy for food availability. & (Geselbracht et al. 2024; Peyre,
Marshall, and Sable 2021) \\
\end{longtable}

\end{landscape}

It should be recognized that not all environmental factors are
adequately spatially distributed or represented within the available
datasets, and therefore careful consideration should be given when
selecting the final dataset used for modeling habitat suitability
(Peyre, Marshall, and Sable 2021; Geselbracht et al. 2024; Christensen
et al. 1997).

\subsection{Species Distribution
Modelling}\label{species-distribution-modelling}

Quantifying the relationships between the distribution of species and
the species-environment interactions have a long history in ecological
research, having been described as far back as the late 17th early 18th
century in which researchers observed biological patterns in terms of
their relationships with geographical and/or environmental gradients
(Rushton, Ormerod, and Kerby 2004; Elith and Leathwick 2009). Since
then, the propagation of large-scale spatial datasets, the development
and integration of GIS, increased access to statistical packages and
computational power have allowed for much more development in modeling
complex species-environment interactions (Johnson and Gillingham 2005).
Out of such developments has arisen a technique known as Species
Distribution Modeling (SDM), which can be loosely defined as numerical
tools that combine observations of species occurrence or abundance with
environmental estimates (Elith and Leathwick 2009). According to Elith
and Leathwick (2009), SDMs can be largely classified within several
families of models: expert-driven or heuristic, environmental-envelope,
environmental-similarity, regression-based, Bayesian and hierarchical,
machine-learning and data mining, and ensemble or multimodal.
Additionally, the selection of an SDM approach largely depends on a
multitude of factors such as: the study's objective, quality and
quantity of available data, ecological characteristics of the species,
spatial and temporal scale, whether input data includes absence data or
presence-only data (Elith and Leathwick 2009). With so many choices of
modeling technique and the driving factors behind that choice, deciding
on one model can be quite difficult and the choice usually defaults to a
standard technique depending on the subject matter.

\emph{Habitat Suitability Indices}

One popular approach used in many studies involving the Eastern oyster
throughout the gulf region is that of the expert-driven Habitat
Suitability Index (Barnes et al. 2007; Pollack et al. 2012; Geselbracht
et al. 2024; Christensen et al. 1997). HSIs are typically used in
situations where quantitative data on species distributions is difficult
to collect or non-existent, relying on expert-driven variable selection,
meta-analysis, and classification methods that transform environmental
characteristics into standardized estimates of habitat suitability
(Barnes et al. 2007; Pollack et al. 2012; Geselbracht et al. 2024;
Johnson and Gillingham 2005).

Within the PBS, HSIs have been used to develop GIS-based and
estuary-specific HSMs that integrate biological, physical, chemical, and
management-related variables in order to identify suitable restoration
areas with the system. In the mid 1990's Christensen et al. (1997)
developed a HSI for the eastern oyster in the PBS to demonstrate the
model frameworks usefulness in illustrating species suitability in a
system wide approach. Geselbracht et al. (2024) then followed and
further expanded on the approach incorporating management factors into
the model, adding a layer of regulatory influence to the model output
versus the purely environmental analysis of the prior. Additionally, the
Pensacola and Perdido Bay Estuary Program recently built on-top of
Geselbracht et al. (2024) by incorporating some of their own data and
refining some of the weighting (Pensacola and Perdido Bays Estuary
Program n.d.).

These previous studies alone highlight the strengths and weakness of the
HSI approaches. Some of the strengths include interoperability,
adaptability, and the integration of ecological knowledge, but the
results entirely rely on expert decisions regarding variable selection,
suitability functions, weighting and scoring, and how to aggregate the
results. All of these present opportunity for differing outputs even if
the same inputs were used by different groups (Geselbracht et al. 2024;
Johnson and Gillingham 2005).

\emph{Maximum Entropy Modelling (MaxEnt)}

MaxEnt is a presence-background SDM framework based on general-purpose
machine learning methods that are useful for making predictions using
incomplete information by relating environmental conditions at known
occurrence locations to the environmental conditions throughout the
study area (Phillips, Anderson, and Schapire 2006; Elith et al. 2011).
Some of the main benefits of MaxEnt is that it uses presence-only data,
reducing the need for the often un-available absence data, it has the
ability to model complex interactions between environmental conditions,
and it has demonstrated high predictive performance while limiting over
fitting through regularization (Phillips et al. 2009; Merow, Smith, and
Silander 2013; Elith et al. 2011). An added benefit to this particular
study is the fact that MaxEnt is tailored towards presence-only species
occurrence data, making it highly suitable for the data set used here in
which the FWC-FWRI (2026) datasets explicitly collected occurrence data
without including absence locations.

However, the use of such algorithms should not be made blindly as many
issues can arise if careful precautions are not taken. Due to the nature
of ecological many biases in the data samples can arise such as spatial
and temporal autocorrelation, group dependence structures, hierarchical
structures and multicollinearity leading to overfitting model parameters
and over-optimistic prediction confidence (Dormann et al. 2013; Kass et
al. 2021; Roberts et al. 2017). Thankfully, significant work has been
put into this through cross-validation techniques and the development of
model evaluation software packages such as ENMeval 2.0 implemented
through the R statistical language (Kass et al. 2021; Roberts et al.
2017)

\textbf{TODO: THINK OF BETTER SECTION NAME} \emph{Hypothesis}

To demonstrate how a MaxEnt SDM can be used as a complimentary tool to
help researchers and practitioners leading ecological restoration
activities, this study aims to develop an empirically fitted MaxEnt
habitat-suitability model using observed oyster occurrence data and
known environmental conditions. The primary objective will be to develop
a spatially explicit SDM representing relative Eastern oyster habitat
suitabiliy throughout the PBS. Next, the resulting suitability surface
will be evaluated agains the existing Pensacola Bay HSI developed by
Geselbracht et al. (2024) and iterated upon by Pensacola and Perdido
Bays Estuary Program (n.d.). To achieve this, environmental predictors
that are commonly associated with oyster habitat suitability will be
identified and obtained. The predictors will be used to fit the MaxEnt
model and model performance will be evaluated using spatially separated
cross-validation. Finally, we will compare the predicted surface to the
existing HSI to determine areas of agreement and disagreement, and we
will use these results to determine whether the MaxEnt framework
provides an empirical alternative or compliment to HSI-based restoration
planning. It is hypothesized that the MaxEnt model will successfully
identify patterns of habitat suitability for the Easter oyster withinthe
PBS. Overall agreement between the HSI and MaxEnt is expected to be
high. However, MaxEnt is expected to produce a smoother, more continuous
suitability surface because of the empirically fitted environmental
relationships as opposed to predefined, lower-resolution HSI suitability
envelopes. As a result, the MaxEnt model could prove beneficial to help
narrow down project locations to help reinforce expert-driven model
outputs.

\section{Methods}\label{methods}

\emph{Study design and team status}

This capstone is being completed as an individual project. As such, all
data acquisition, processing, modeling, and interpretations will be
produced soley by the author.

All data analysis and processing will be performed using the R
programming language (R Core Team 2024). Environmental observation data
will be acquired from the EPA's Water Quality Portal via the
dataRetrieval r-package (USGS 2026; DeCicco et al. 2026). Data QA/QC and
preprocessing will be completed utilizing the Tidyverse package and any
spatially aware data processing will be conducted using the package sf
(Wickham et al. 2019; Pebesma 2018). Environmental data will be
processed to create continuous surface rasters for model input using the

\emph{Study Area}

The study area for this project is the Pensacola Bay System which
includes portions of Pensacola Bay, Escambia Bay, East Bay and
Blackwater Bay. Boundaries of this study area were achieved by accessing
Florida Department of Protections (FDEP) Waterbod IDs feature layer
hosted on ArcGIS Online (FDEP 2005). Once retrieved, the dataset was
filtered to include only relevant WBIDs (WBID = 548) and then combined
to create a singular study area polygon. Creating the study area polygon
is an important first step as it will be needed to extract relevant
environmental data from much larger datasets in the coming data
gathering. The study area was intentionally chosen in an attempt to
match the previous study area applied in Geselbracht et al. (2024).

\begin{figure}

\centering{

\pandocbounded{\includegraphics[keepaspectratio]{4-tables-and-figures/study-area.png}}

}

\caption{\label{fig-study-area}Study area for the Pensacola Bay System.}

\end{figure}%

\emph{Data acquisition and dataset inventory}

Oyster occurence data was acquired from the Florida Fish and Wildlife
Commission's GIS portal as polygons and converted to points using
polygon centroids. After clipping to the study area extent, \#\#\# known
reef locations were present within the PBS (FWC-FWRI 2026). Salinity,
dissolved oxygen, temperature, and turbidity, and chlorophyll-a
observations were obtained from the Water Quality Portal using the
`dataRetrieval' package in R for all sites located within the PBS and
collected within the 20 year period from 2005 to 2025 DeCicco et al.
(2026). Parameters were chosen based on relevance and represenation in
previous oyster habitat retoration studies (Peyre, Marshall, and Sable
2021; Barnes et al. 2007; Pollack et al. 2012; Geselbracht et al. 2024).
Bathymetry data for the PBS was obtained from NOAA as a \#\#\#\#
resolution raster dataset (National Centers for Environmental
Information 2018). \textbf{\emph{DOES BATHYMETRY DATA NEED TO BE
RESAMPLED TO NEW CELL SIZE?}}. Substrate data was obtained from the
Pensacola and Perdido Bay Estuary Program and was the same dataset used
in their earlier HSI study which was modeled on the
(\textbf{Geselbrach2024?}) study.


\begin{longtable}[]{@{}
  >{\raggedright\arraybackslash}p{(\linewidth - 8\tabcolsep) * \real{0.2235}}
  >{\raggedright\arraybackslash}p{(\linewidth - 8\tabcolsep) * \real{0.1176}}
  >{\raggedleft\arraybackslash}p{(\linewidth - 8\tabcolsep) * \real{0.1529}}
  >{\raggedleft\arraybackslash}p{(\linewidth - 8\tabcolsep) * \real{0.1882}}
  >{\raggedright\arraybackslash}p{(\linewidth - 8\tabcolsep) * \real{0.3176}}@{}}

\caption{\label{tbl-metadata}Environmental Parameters}

\tabularnewline

\toprule\noalign{}
\begin{minipage}[b]{\linewidth}\raggedright
Dataset
\end{minipage} & \begin{minipage}[b]{\linewidth}\raggedright
Data Type
\end{minipage} & \begin{minipage}[b]{\linewidth}\raggedleft
\# of Records
\end{minipage} & \begin{minipage}[b]{\linewidth}\raggedleft
Processed Count
\end{minipage} & \begin{minipage}[b]{\linewidth}\raggedright
Source
\end{minipage} \\
\midrule\noalign{}
\endhead
\bottomrule\noalign{}
\endlastfoot
Salinity & Point & 6992 & 95 &
(USGS 2026) \\
Water Temperature & Point & 7159 & 100 &
(USGS 2026) \\
Substrate type & Raster & NA & NA &
(Pensacola and Perdido Bays Estuary Program n.d.) \\
Dissolved oxygen & Point & 6261 & 99 &
(USGS 2026) \\
Depth / bathymetry & Raster & NA & NA &
(USGS 2026) \\
Turbidity & Point & 3402 & 75 &
(USGS 2026) \\
Chlorophyll a & Point & 1847 & 44 &
(USGS 2026) \\
Oyster Reefs & Polygon & 443 & 443 &
(FWC-FWRI 2026) \\

\end{longtable}

\textbf{Water-quality Data Preprocessing}

All water quality parameters were inspected to ensure comparable
measurement units and analytical procedures and avoid combining
different forms of similar analytes. Further, to maintain consistency
with other similar oyster studies, all data was filtered to only include
warm season conditions. Here, the warm season is defined as May through
October. Finally, the data was screened for any extreme outliers, where
outliers were defined as: \[
\begin{aligned}
\text{Outlier if } \quad
x &< Q_{0.25} - 1.5(IQR) \\
\text{or} \quad
x &> Q_{0.75} + 1.5(IQR) \\
\\
\text{where } IQR &= Q_{0.75} - Q_{0.25}
\end{aligned}
\]

Following data cleaning, station-level environmental summaries were
calculated to represent warm-season conditions that could be expected at
that station. Stations that only had one sample were also removed as it
was considered one sample was not representitive enough to be considered
the normal conditions over the full 20 year period. Finally, the
station-level dataset was converted to spatial using the \texttt{sf}
package in r and projected to NAD 1983 (2011) StatePlane Florida North
FIPS 0903 in meters (Pebesma 2018).

\textbf{Spatial Interpolation of Environmental Factors}

Spatial interpolation in progress\ldots.

\section{Results}\label{results}

\emph{Water Quality Observations}

Document water-quality dataset size and temporal coverage (USGS 2026). -
Dissolved oxygen: approximately 145 stations and 3,072 observations . -
Salinity: approximately 135 stations and 3,401 observations -
Temperature: approximately 145 stations and 3,442 observations -
Turbidity: approximately 122 stations and 1,563 observations -
Chlorophyll-a: approximately 48 stations and 754 observations - Current
temporal coverage is approximately 2005--2024

\section{Discussion}\label{discussion}

In progress\ldots{}

\section{Conclusion}\label{conclusion}

In progress\ldots{}

\section{References}\label{references}

\phantomsection\label{refs}
\begin{CSLReferences}{1}{0}
\bibitem[\citeproctext]{ref-Barnes2007}
Barnes, T. K., A. K. Volety, K. Chartier, F. J. Mazzotti, and L.
Pearlstine. 2007. {``A Habitat Suitability Index Model for the Eastern
Oyster ({Crassostrea virginica}), a Tool for Restoration of the
Caloosahatchee Estuary, Florida.''} \emph{Journal of Shellfish Research}
26 (4): 949--59.
\url{https://doi.org/10.2983/0730-8000(2007)26\%5B949:AHSIMF\%5D2.0.CO;2}.

\bibitem[\citeproctext]{ref-Birch2021}
Birch, A., R. Brumbaugh, B. DeAngelis, L. Geselbracht, A. Graves, J.
Blair, and R. Jones. 2021. {``Oyster Fisheries and Habitat Management
Plan for the Pensacola Bay System.''} The Nature Conservancy.
\url{https://www.nature.org/content/dam/tnc/nature/en/documents/PBS_FloridaOyster-Fisheries-Habitat-Mgt-Plan_18May2021_Final.pdf}.

\bibitem[\citeproctext]{ref-Christensen1997}
Christensen, J D, T A Battista, M E Monaco, and C J Klein. 1997.
{``Habitat Suitability Index Modeling and GIS Technology to Support
Habitat Management: Pensacola Bay, Florida Case Study.''} U.S. National
Oceanic; Atmospheric Administration.
\url{https://repository.library.noaa.gov/view/noaa/55683}.

\bibitem[\citeproctext]{ref-dataRetrieval2026}
DeCicco, Laura, Robert Hirsch, David Lorenz, Jordan Read, Jordan Walker,
Lindsay Platt, David Watkins, et al. 2026. {``dataRetrieval: R Packages
for Discovering and Retrieving Water Data Available from u.s. Federal
Hydrologic Web Services.''} Reston, VA: U.S. Geological Survey; U.S.
Geological Survey. \url{https://doi.org/10.5066/P9X4L3GE}.

\bibitem[\citeproctext]{ref-Dormann2013}
Dormann, Carsten F., Jane Elith, Sven Bacher, Carsten Buchmann, Gudrun
Carl, Gabriel Carré, Jaime R.García Marquéz, et al. 2013.
{``Collinearity: A Review of Methods to Deal with It and a Simulation
Study Evaluating Their Performance.''} \emph{Ecography} 36: 27--46.
\url{https://doi.org/10.1111/j.1600-0587.2012.07348.x}.

\bibitem[\citeproctext]{ref-Elith2009}
Elith, Jane, and John R. Leathwick. 2009. {``Species Distribution
Models: Ecological Explanation and Prediction Across Space and Time.''}
\emph{Annual Review of Ecology, Evolution, and Systematics} 40
(December): 677--97.
\url{https://doi.org/10.1146/annurev.ecolsys.110308.120159}.

\bibitem[\citeproctext]{ref-Elith2011}
Elith, Jane, Steven J. Phillips, Trevor Hastie, Miroslav Dudík, Yung En
Chee, and Colin J. Yates. 2011. {``A Statistical Explanation of MaxEnt
for Ecologists.''} \emph{Diversity and Distributions} 17 (January):
43--57. \url{https://doi.org/10.1111/j.1472-4642.2010.00725.x}.

\bibitem[\citeproctext]{ref-FDEP2005WBIDs}
FDEP. 2005. {``Waterbody IDs (WBIDs).''} Geospatial Open Data; Florida
Department of Environmental Protection.
\url{https://geodata.dep.state.fl.us/datasets/FDEP::waterbody-ids-wbids/about}.

\bibitem[\citeproctext]{ref-FWC2026Oysters}
FWC-FWRI. 2026. {``Oyster Beds in Florida.''} Geospatial Open Data;
Florida Fish; Wildlife Conservation Commission, Fish; Wildlife Research
Institute.
\url{https://geodata.myfwc.com/datasets/myfwc::oyster-beds-in-florida/about}.

\bibitem[\citeproctext]{ref-Geselbracht2024}
Geselbracht, L., M. Johnston, B. M. DeAngelis, and A. Birch. 2024.
{``Estuary-Specific and Adaptive Habitat Suitability Index Model for the
Eastern Oyster Crassostrea Virginica in the Pensacola Bay System,
Florida, USA.''} \emph{Coastal Management} 52: 17--34.
\url{https://doi.org/10.1080/08920753.2024.2335129}.

\bibitem[\citeproctext]{ref-Grabowski2012}
Grabowski, Jonathan H., Robert D. Brumbaugh, Robert F. Conrad, Andrew G.
Keeler, James J. Opaluch, Charles H. Peterson, Michael F. Piehler, Sean
P. Powers, and Ashley R. Smyth. 2012. {``Economic Valuation of Ecosystem
Services Provided by Oyster Reefs.''} \emph{BioScience} 62 (October):
900--909. \url{https://doi.org/10.1525/bio.2012.62.10.10}.

\bibitem[\citeproctext]{ref-Johnson2005}
Johnson, Chris J., and Michael P. Gillingham. 2005. {``An Evaluation of
Mapped Species Distribution Models Used for Conservation Planning.''}
\emph{Environmental Conservation} 32 (June): 117--28.
\url{https://doi.org/10.1017/S0376892905002171}.

\bibitem[\citeproctext]{ref-Kass2021}
Kass, Jamie M., Robert Muscarella, Peter J. Galante, Corentin L. Bohl,
Gonzalo E. Pinilla-Buitrago, Robert A. Boria, Mariano Soley-Guardia, and
Robert P. Anderson. 2021. {``ENMeval 2.0: Redesigned for Customizable
and Reproducible Modeling of Species' Niches and Distributions.''}
\emph{Methods in Ecology and Evolution} 12 (September): 1602--8.
\url{https://doi.org/10.1111/2041-210X.13628}.

\bibitem[\citeproctext]{ref-Merow2013}
Merow, Cory, Matthew J. Smith, and John A. Silander. 2013. {``A
Practical Guide to MaxEnt for Modeling Species' Distributions: What It
Does, and Why Inputs and Settings Matter.''} \emph{Ecography} 36
(October): 1058--69.
\url{https://doi.org/10.1111/j.1600-0587.2013.07872.x}.

\bibitem[\citeproctext]{ref-NOAA_2018}
National Centers for Environmental Information. 2018. {``NOAA NOS
Estuarine Bathymetry - Pensacola Bay (G130).''}
\url{https://doi.org/10.7289/V5XK8CWK}.

\bibitem[\citeproctext]{ref-sf}
Pebesma, Edzer. 2018. {``{Simple Features for R: Standardized Support
for Spatial Vector Data}.''} \emph{{The R Journal}} 10 (1): 439--46.
\url{https://doi.org/10.32614/RJ-2018-009}.

\bibitem[\citeproctext]{ref-PPBEP_OysterRestoration}
Pensacola and Perdido Bays Estuary Program. n.d. {``Pensacola Bay Oyster
Restoration Initiative.''}
\url{https://www.ppbep.org/what-we-do/restoration/pensacola-bay-oyster-restoration-initiative}.

\bibitem[\citeproctext]{ref-LaPeyre2021}
Peyre, Megan K. La, Danielle A. Marshall, and Shaye E. Sable. 2021.
{``Oyster Model Inventory: Identifying Critical Data and Modeling
Approaches to Support Restoration of Oyster Reefs in Coastal u.s. Gulf
of Mexico Waters.''} \emph{Diversity and Distributions}. Vol. 17. U.S.
Geological Survey. \url{https://doi.org/10.3133/ofr20211063}.

\bibitem[\citeproctext]{ref-Phillips2006}
Phillips, Steven J., Robert P. Anderson, and Robert E. Schapire. 2006.
{``Maximum Entropy Modeling of Species Geographic Distributions.''}
\emph{Ecological Modelling} 190 (January): 231--59.
\url{https://doi.org/10.1016/j.ecolmodel.2005.03.026}.

\bibitem[\citeproctext]{ref-Phillips2009}
Phillips, Steven J., Miroslav Dudík, Jane Elith, Catherine H. Graham,
Anthony Lehmann, John Leathwick, and Simon Ferrier. 2009. {``Sample
Selection Bias and Presence-Only Distribution Models: Implications for
Background and Pseudo-Absence Data.''} \emph{Ecological Applications} 19
(January): 181--97. \url{https://doi.org/10.1890/07-2153.1}.

\bibitem[\citeproctext]{ref-Pollack2012}
Pollack, Jennifer Beseres, Andrew Cleveland, Terence A. Palmer, Anthony
S. Reisinger, and Paul A. Montagna. 2012. {``A Restoration Suitability
Index Model for the Eastern Oyster (Crassostrea Virginica) in the
Mission-Aransas Estuary, TX, USA.''} \emph{PLoS ONE} 7 (July).
\url{https://doi.org/10.1371/journal.pone.0040839}.

\bibitem[\citeproctext]{ref-r}
R Core Team. 2024. \emph{R: A Language and Environment for Statistical
Computing}. Vienna, Austria: R Foundation for Statistical Computing.
\url{https://www.R-project.org/}.

\bibitem[\citeproctext]{ref-Roberts2017}
Roberts, David R., Volker Bahn, Simone Ciuti, Mark S. Boyce, Jane Elith,
Gurutzeta Guillera-Arroita, Severin Hauenstein, et al. 2017.
{``Cross-Validation Strategies for Data with Temporal, Spatial,
Hierarchical, or Phylogenetic Structure.''} \emph{Ecography}. Blackwell
Publishing Ltd. \url{https://doi.org/10.1111/ecog.02881}.

\bibitem[\citeproctext]{ref-Rushton2004}
Rushton, S. P., S. J. Ormerod, and G. Kerby. 2004. {``New Paradigms for
Modelling Species Distributions?''} \emph{Journal of Applied Ecology}.
\url{https://doi.org/10.1111/j.0021-8901.2004.00903.x}.

\bibitem[\citeproctext]{ref-TNC2021OysterRestoration}
The Nature Conservancy. 2021. {``Pensacola East Bay Oyster Habitat
Restoration Project.''} The Nature Conservancy.
\url{https://www.nature.org/content/dam/tnc/nature/en/documents/PEB-Oyster-Habitat-Restoration-Fact-Sheet-FINAL-7-13-21.pdf}.

\bibitem[\citeproctext]{ref-wqp_2026}
USGS, USDA, USEPA. 2026. {``Water Quality Portal.''} {U.S. Geological
Survey}; {U.S. Environmental Protection Agency}; {U.S. Department of
Agriculture}. \url{https://www.waterqualitydata.us/}.

\bibitem[\citeproctext]{ref-tidyverse}
Wickham, Hadley, Mara Averick, Jennifer Bryan, Winston Chang, Lucy
D'Agostino McGowan, Romain François, Garrett Grolemund, et al. 2019.
{``Welcome to the {tidyverse}.''} \emph{Journal of Open Source Software}
4 (43): 1686. \url{https://doi.org/10.21105/joss.01686}.

\end{CSLReferences}




\end{document}
